% ---------------------------------------------------------------------------
\section{The voynichizer}
\label{sec:generator}

The voynichizer \citep{voynichizer} takes any text in UTF-8 and a key, and writes a book with the 207 folio names,
sections and Currier languages of the manuscript, about 4{,}050--4{,}260 lines in EVA transliteration. With the key the
text comes back exactly; with a wrong key it is rejected. This section describes version~21, the published one.
Figure~\ref{fig:pipeline} gives the overall design and Figure~\ref{fig:page} an example.

\begin{figure}[htbp]
\centering
\begin{tikzpicture}[font=\footnotesize,
  box/.style={draw, rounded corners=2pt, align=center, inner sep=3pt, minimum height=10mm, text width=#1},
  box/.default=22mm,
  key/.style={draw, dashed, rounded corners=2pt, align=center, inner sep=3pt, minimum height=7mm, text width=14mm},
  arr/.style={-{Stealth[length=2mm]}, thick}]
\node[box=16mm] (msg) at (0,0) {text\\(any language)};
\node[box=24mm] (enc) at (2.85,0) {compress, frame, encrypt, authenticate};
\node[box=17mm] (dec) at (5.85,0) {arithmetic decoder};
\node[box=25mm] (cnt) at (8.95,0) {how many times each known word appears on each page};
\node[box=15mm] (book) at (12.05,0) {book of 207 folios, EVA};
\node[key] (k) at (2.85,2.0) {key};
\node[box=46mm] (model) at (7.4,2.0) {model of the manuscript: lexicon of each section and language, character of a
real page, invented words};
\node[box=26mm] (arr) at (8.95,-2.0) {layout and word order (no information)};
\node[box=26mm] (spec) at (3.6,-2.0) {scorecard, line gate, judges (Section~\ref{sec:scorecard})};
\draw[arr] (msg) -- (enc);
\draw[arr] (enc) -- node[above]{bits} (dec);
\draw[arr] (dec) -- (cnt);
\draw[arr] (cnt) -- (book);
\draw[arr] (k) -- (enc);
\draw[arr] (k) -- (model);
\draw[arr] (model.south -| dec.north) -- node[left]{probabilities} (dec.north);
\draw[arr] (arr) -| (book);
\draw[arr, dashed] (spec) -- node[above]{tuned against} (arr);
\end{tikzpicture}
\caption{Design of the voynichizer. The encrypted bits only decide the word counts of each page, through an arithmetic
decoder that samples them from the model; everything else comes from the key and the manuscript's statistics. Reading
runs the arithmetic encoder on the observed counts.}
\label{fig:pipeline}
\end{figure}

\subsection{Design principle}

Feeding uniformly random bits to an arithmetic decoder produces an exact sample of the decoder's probability model, up to
the quantisation of the probabilities. If the bits are the ciphertext of a message under a good cipher, they are
indistinguishable from random bits, and the book is a sample of the model whether or not it carries a message: security
relative to the model in the sense of \citet{cachin1998}, by the method of \citet{ziegler2019}. Two consequences follow.
The model, not the message, decides how close the book comes to the manuscript; and the capacity is fixed by the entropy
of the part of the model that carries the bits, and cannot be raised without changing the statistics.

\subsection{From text to bits}

The text is compressed with zlib and framed as \emph{length} (4 bytes) $\|$ \emph{tag} (8 bytes) $\|$
\emph{compressed text}, where the tag is a truncated HMAC-SHA256. The frame is XORed with a SHAKE-256 keystream; after
it, the same keystream continues as filler, so that a book without a message is driven by the keystream alone. The
master key is derived with scrypt ($n = 2^{15}$, $r = 8$, $p = 1$) from the passphrase; each random choice of the
program has its own generator seeded from an HMAC of the master key and the use (and the page, for the choices made
page by page). For example, the opening of
Book~XVII of Isidore's \emph{Etymologiae} (11{,}207 bytes) compresses to 4{,}776 bytes and becomes 38{,}304 framed bits.

\subsection{The model of the book}

\paragraph{Layout} For each page type (section and Currier language), the program keeps the number of lines per page, the
page width (median words per line) and paragraph lengths, and for the whole book the ratio of words in a line to the
page width by line role (first, middle, last, only line of a paragraph). It draws a page from these statistics, so no
generated page copies the layout of a real one; books come out about 5\% longer than the manuscript.

\paragraph{Known and invented words} About 86\% of the manuscript's running words are \emph{known} words (2{,}246 types
that occur at least twice); the rest occur once. In each of eight place classes (first line of a paragraph or not, times
first, second, middle or last word of the line) a place receives an invented word with the probability observed for
once-occurring words in that class in the manuscript, adjusted by the page. Invented words are generated by a glyph model
(each glyph predicted from the three before it, backing off to two) learned on the 4{,}776 once-occurring words, with
lengths drawn by place class; an invented word of six glyphs or more is, with probability 0.25, replaced if possible by
a join of two known words of the page whose junction is probable. An invented word is never a Voynich word,
so it cannot be confused with a known one when reading.

\paragraph{The lexicon of a page and its ``character''} The known-word lexicon of a page is made of the occurrences of
known words on the other pages of the same section and language (never the page itself). The key chooses another real
page of the same type, whose glyph frequencies give the page its ``character'': each word's weight is its count raised to
$\gamma = 0.976$, times $\exp(\kappa \sum_g n_g(w)\,\delta_g)$, where $\delta_g$ is the smoothed log-ratio of glyph $g$
between the chosen page and its section and $\kappa = 0.644$.

\subsection{The channel: word counts}

The bits decide only how many times each known word appears on each page. For a page with $n$ known-word places and
ordered integer weights $w_1 \ge \dots \ge w_m$ (total $W$), the counts are one multinomial draw, written as a chain of
binomials: $c_1 \sim \mathrm{Bin}(n, w_1/W)$, $c_2 \sim \mathrm{Bin}(n - c_1, w_2/(W - w_1))$, and so on. Each binomial
is a run of binary decisions (``stop at $j$ copies?'', with probability $P(X = j)/P(X \ge j)$ quantised to 16 bits), and
each decision is a symbol of a 32-bit binary arithmetic coder. Writing runs the coder's \emph{decoder} on the ciphertext
bits, so the decisions are sampled; reading runs its \emph{encoder} on the observed decisions and recovers the bits, then
XORs them, checks length and tag and decompresses. The bits are consumed in page order; the Isidore text fills about 130
of the 207 pages and the filler the rest.

\paragraph{Capacity} A v21 book carries 81{,}100--85{,}900 bits (mean 83{,}400 over 24 keys), about 2.7 bits per known
word, roughly 20{,}000 characters of ordinary prose after compression. The first versions hid the bits in the five
spelling choices (\eva{ch}/\eva{sh}, \eva{k}/\eva{t}, \eva{-l}/\eva{-r}, \eva{qo-}/\eva{o-}, \eva{-dy}/\eva{-ey}), with
about 0.8 bits per choice and 46{,}000 bits per book; moving the channel into the word counts nearly doubled the capacity
and made the books harder to tell apart (Table~\ref{tab:versions}).

\subsection{Arrangement: word order carries no information}

Once the multiset of words of a page is fixed, its order is chosen by a Metropolis sampler (60 proposed swaps per word,
temperature 1) on a score with thirteen terms learned or tuned on the manuscript: the affinity of each word for its place
class (a logistic model on the word's first and last glyphs, length, prefix and ending); links between neighbouring words
(ending--prefix and prefix--prefix, last glyph--first glyph, joins that form a word of the book, pairs seen on other
pages, identical neighbours); similarity with the word at the same position in the line above and with the word two
places away; avoidance of a line starting with the same glyph as the line above; agreement of the five spelling choices
within a line and between consecutive lines; a difference between the two halves of the page; and line widths. None of
these choices affects reading: moving words within a page changes nothing; adding, removing or changing a known word breaks
the reading, and there is no error correction.

\subsection{Security, as far as it goes}

The building blocks are standard, but the whole has not been reviewed by a cryptographer. The scrypt salt is fixed, so a
dictionary of guessed passphrases can be prepared once for all books; the keystream depends only on the key, so two books
written with the same key share it and a different key should be used for each message. A long random passphrase (six
random words or twelve random characters) is needed. The program hides \emph{that there is a message}, not that the
program was used: about 85\% of the running words of a book are Voynich words, and a book is recognisable as the
program's output.

\subsection{Development}

Table~\ref{tab:versions} summarises the versions that changed the method. Every version was judged on the same
instruments. From version~14 on, the numbers we cite come from keys never used for tuning; versions 17 and 20 were
released on the development keys and confirmed on the other 12 later the same day, version~21 before release.

\begin{table}[htbp]
\centering\footnotesize
\caption{Main versions of the voynichizer (AUC of the two judges; the early v8 benchmark used one key with three seeds
and is optimistic compared with later replications).}
\label{tab:versions}
\begin{tabularx}{\linewidth}{lYr}
\toprule
version & what changed & AUC judge 1 / judge 2 \\
\midrule
v0 & body from an earlier generator; message forced into the five spelling choices & 0.958 / 0.980 \\
v1 & arithmetic coding over a model of the spelling choices; about 46{,}000 bits & 0.889 / 0.957 \\
v8 & generator ``by pieces'': page lexicon, borrowed character, invented words, learned arrangement & 0.604 / 0.715 (3 seeds) \\
v10 & message moved into the word counts & 0.572 / 0.690 (12 keys) \\
v14 & joins accepted only with a probable junction & 0.575 / 0.600 (24 keys) \\
v15--v17 & authenticated encryption (scrypt, SHAKE-256, HMAC); statistical layout; first public release & 0.559 / 0.597 (24 keys) \\
v21 & correction of an initialisation defect in the arrangement (found by reading the public code) & 0.554 / 0.601 (24 keys) \\
\bottomrule
\end{tabularx}
\end{table}

\begin{figure}[p]
\centering\footnotesize
% generato da genera_figure.py: prime 12 righe di f1r (ZL) e del libro d'esempio v21
\begin{minipage}{\linewidth}
\textbf{f1r, the manuscript} (first 12 lines, ZL transliteration)\par\medskip
{\voyfont\large\raggedright\setlength{\parskip}{2pt}
fachys ykal ar ataiin shol shory cthres y kor sholdy\par
sory ckhar or y kair chtaiin shar ase cthar cthar dan\par
syaiir sheky or ykaiin shod cthoary cthes daraiin sy\par
soiin oteey oteos roloty cthiar daiin okaiin or okan\par
sair y chear cthaiin cphar cfhaiin\par
ydaraishy\par
odar cy shol cphoy oydar sh s cfhoaiin shodary\par
yshey shody okchoy otchol chocthy os chy dain chor kos\par
daiin shos cfhol shody\par
dain os teody\par
ydain cphesaiin ols cphey ytain shoshy cphodal es\par
oksho kshoy otairin oteol okan shodain sckhey daiin
\par}
\vspace{8mm}
\textbf{f1r, a book generated by version 21} (first 12 lines)\par\medskip
{\voyfont\large\raggedright\setlength{\parskip}{2pt}
ocheoporaiin shdar opchey otar cphedy opchedy aiir shkechom\par
soin ain chees oaly oiies ykdain dam chekchy daiin\par
sain olchs ar shedy oltal ched air olchdy okaiin olkeol ytor\par
ytaiin ykchedy chol tor chekain olkal ain okeedy otaiin\par
yrool ytchy oraiin airodeey lchey otain otar\par
olky al otchol okaiin chdal r qo keeo r shckhy olkey qokain or alchdy qokal okal ry\par
sain sheedy qokal shedy lkeear olkeedy qoeekeey qoky\par
ytchy kchedy dlar cheoky lol chekain chor dcheedy chaiin chedy yteol qoky soraiin dar\par
qo lchedy qokey qokaiin oteey lkeol qokaldy ain yokecheeg\par
or aiiin okaiin orkoldy\par
tody keody cheol shedy opchsd opo oteey otedo qokar oteody yfcheety\par
sain olain okain choepy shey ol okeedal cholkaly otain chey ar shedy aly
\par}
\end{minipage}

\caption{The first twelve lines of folio 1r in the manuscript (top) and in a book generated by version~21 (bottom), both
set in the voynichizer's font. The generated book hides the sentence ``This book was written by the voynichizer. Its words hide
this sentence, compressed and encrypted; without the key it reads like the Voynich manuscript and says nothing.'' It can
be read back with the public program and the key \texttt{white paper figure one, public example key}. The message uses
1{,}144 of the book's 83{,}480 bits.}
\label{fig:page}
\end{figure}
